\begin{table}[t]
\centering
\begin{threeparttable}
\caption{Returns at stake per filing-year cohort, by lever and owner group (millions of 2024 PPP dollars)}\label{tab:levers}
\small
% Generated by em-patent-value code/38_paper_exhibits.py (policy_levers.csv, step 26); do not edit by hand
\begin{tabular}{lccccc}
\toprule
 & \multicolumn{3}{c}{Fee relief} & Markets for technology & University patenting \\
\cmidrule(lr){2-4}\cmidrule(lr){5-5}\cmidrule(lr){6-6}
 & China & Korea & Brazil & China & China \\
\midrule
Small firms (1--9)          & 0.07 & 0.49 & 0.05 & 133 &  \\
Mid-sized firms (10--99)    & 0.02 & 0.42 & 0.03 & 146 &  \\
Large firms (100+)          & n.e. & n.e. & n.e. & 0.00 &  \\
Universities                & 0.43 & 0.34 & 0.04 & 462 & 178--319 \\
\addlinespace
All owners                  & 0.52 & 1.3 & 0.11 & 742 & 178--319 \\
\quad \% of national R\&D & 0.00 & 0.00 & 0.00 & 0.79 & 0.19--0.34 \\
\bottomrule
\end{tabular}
\begin{tablenotes}\footnotesize
\item \emph{Notes:} Change in the returns domestic owners earn from their patents, excluding fees, for one filing-year cohort; deterministic renewal model, each patent simulated on its own fee path. Markets for technology: upper bounds, calibrated to the extra post-sale survival of arm's-length sales; the firm targets are imprecise (\SuppTab{S5}{tab:sales}). The two university entries are alternative routes to the same value and are not additive. n.e.: not eligible. Details in footnote~\ref{fn:levers}.
\end{tablenotes}
\end{threeparttable}
\end{table}
